%-------------------------
% Resume in Latex
% Author : Jake Gutierrez
% Based off of: https://github.com/sb2nov/resume
% License : MIT
%------------------------

\documentclass[letterpaper,11pt]{article}

\usepackage{latexsym}
\usepackage[empty]{fullpage}
\usepackage{titlesec}
\usepackage{marvosym}
\usepackage[usenames,dvipsnames]{color}
\usepackage{verbatim}
\usepackage{enumitem}
\usepackage[hidelinks]{hyperref}
\usepackage{fancyhdr}
\usepackage[english]{babel}
\usepackage{tabularx}
\usepackage[normalem]{ulem}
\input{glyphtounicode}



%----------FONT OPTIONS----------
% sans-serif
% \usepackage[sfdefault]{FiraSans}
% \usepackage[sfdefault]{roboto}
% \usepackage[sfdefault]{noto-sans}
% \usepackage[default]{sourcesanspro}

% serif
% \usepackage{CormorantGaramond}
% \usepackage{charter}


\pagestyle{fancy}
\fancyhf{} % clear all header and footer fields
\fancyfoot{}
\renewcommand{\headrulewidth}{0pt}
\renewcommand{\footrulewidth}{0pt}

% Adjust margins
\addtolength{\oddsidemargin}{-0.5in}
\addtolength{\evensidemargin}{-0.7in}
\addtolength{\textwidth}{1in}
\addtolength{\topmargin}{-0.75in}
\addtolength{\textheight}{1.0in}


\urlstyle{same}

%\raggedbottom
\raggedright
\setlength{\tabcolsep}{0in}

% Sections formatting
\titleformat{\section}{
  \vspace{-10pt}\scshape\raggedright\large
}{}{0em}{}[\color{black}\titlerule \vspace{-5pt}]

% Ensure that generate pdf is machine readable/ATS parsable
\pdfgentounicode=1

%-------------------------
% Custom commands
% Adjust underline depth to be closer to text
\renewcommand{\ULdepth}{2.5pt}

\newcommand{\resumeItem}[1]{
  \item\small{
    {#1 \vspace{-3pt}}
  }
}

\newcommand{\resumeSubheading}[4]{
  \vspace{-2pt}\item
    \begin{tabular*}{0.97\textwidth}[t]{l@{\extracolsep{\fill}}r}
      \textbf{#1} & #2 \\
      \textit{\small#3} & \textit{\small #4} \\
    \end{tabular*}\vspace{-7pt}
}

\newcommand{\resumeSubSubheading}[2]{
    \vspace{-7pt}\item
    \begin{tabular*}{0.97\textwidth}{l@{\extracolsep{\fill}}r}
      \textit{\small#1} & \textit{\small #2} \\
    \end{tabular*}\vspace{-7pt}
}

\newcommand{\resumeProjectHeading}[2]{
    \item
    \begin{tabular*}{0.97\textwidth}{l@{\extracolsep{\fill}}r}
      \small#1 & #2 \\
    \end{tabular*}\vspace{-7pt}
}

\newcommand{\resumeSubItem}[1]{\resumeItem{#1}\vspace{-4pt}}

\renewcommand\labelitemii{$\vcenter{\hbox{\tiny$\bullet$}}$}

\newcommand{\resumeSubHeadingListStart}{\begin{itemize}[leftmargin=0.15in, label={}]}
\newcommand{\resumeSubHeadingListEnd}{\end{itemize}}
\newcommand{\resumeItemListStart}{\begin{itemize}}
\newcommand{\resumeItemListEnd}{\end{itemize}\vspace{-5pt}}

%-------------------------------------------
%%%%%%  RESUME STARTS HERE  %%%%%%%%%%%%%%%%%%%%%%%%%%%%


\begin{document}

%----------HEADING----------
\begin{center}
    \textbf{\Huge \scshape Tony Yiding Tian} \\ \vspace{1pt}
    \small Philadelphia, PA $|$ (267)249-1202 $|$ \href{mailto:tonytg@engineering.upenn.edu}{\underline{tonytg@engineering.upenn.edu}} $|$
    \href{https://tonyxtian.com}{\underline{tonyxtian.com}} $|$
    \href{https://github.com/tonytgrt}{\underline{github.com/tonytgrt}}
\end{center}

%-----------EDUCATION-----------
\section{Education}
  \resumeSubHeadingListStart
    \resumeSubheading
      {University of Pennsylvania}{Philadelphia, PA}
      {B.S.E. in Computer Engineering \& Accelerated M.S.E. in Computer Graphics}{May 2027}
      \resumeItemListStart
        \resumeItem{GPA: 3.79 | Relevant Courses: Operating Systems Implementation, Data Structures \& Algorithms, Computer Organization \& Design, GPU Programming, Embedded Systems, iOS Programming, Advanced Rendering, Machine Learning}
      \resumeItemListEnd
  \resumeSubHeadingListEnd

%-----------EXPERIENCE-----------
\section{Experience}
  \resumeSubHeadingListStart
    \resumeSubheading
      {UI Rendering GPU Software Engineer Intern}{May 2026 -- Aug 2026}
      {Apple}{Cupertino, CA}
      \resumeItemListStart
        \resumeItem{Designed and shipped internal developer tooling for Apple's Core Animation GPU rendering stack, from proposal through design, implementation, on-device validation, and multi-round code review with senior engineers; merged into the main branch and sped up a core internal development workflow by 4x}
        \resumeItem{Engineered a thread-safe publish/observe mechanism in a large, performance-sensitive C++ codebase, replacing a per-consumer duplication scheme so a shared GPU resource is built once and referenced by every consumer}
        \resumeItem{Applied compiler-enforced thread-safety annotations to shared mutable state so lock violations fail the build, and added runtime assertions for an invariant static analysis could not express across differing ownership models}
        \resumeItem{Reduced net code size while adding functionality by unifying two near-duplicate cache-invalidation paths behind one parameterized helper, eliminating unnecessary GPU pipeline recompilation on the narrower path}
        \resumeItem{Triaged and fixed system UI rendering bugs throughout the iOS/macOS 27 beta cycle, diagnosing defects reproducible only in specific SDK and CPU-architecture build configurations; built cross-platform CLI tooling and wrote architecture docs verified against source}
    \resumeItemListEnd

    \resumeSubheading
      {Linux Kernel Policies Research Assistant - PURM Scholar}{May 2025 -- Aug 2025}
      {Learning Directed Operating System (\href{https://ldos.utexas.edu/}{\uline{LDOS}}), Prof. Sebastian Angel}{Philadelphia, PA}
      \resumeItemListStart
        \resumeItem{Developed eBPF-based kernel monitoring infrastructure collecting real-time TCP metrics from 5 core kernel functions (tcp\_v4\_rcv, tcp\_v4\_connect, tcp\_rcv\_state\_process, congestion control, and CUBIC)}
        \resumeItem{Engineered a data analysis pipeline processing 10,000+ TCP state transitions per second to identify performance bottlenecks in kernel networking policies and congestion control algorithms}
        \resumeItem{Contributed 2,000+ lines of C and Python to the open-source \href{https://github.com/tonytgrt/KernMLOps}{\uline{KernMLOps}} repo, implementing kernel probing infrastructure used by 15+ researchers}
    \resumeItemListEnd
  \resumeSubHeadingListEnd

%-----------PROJECTS-----------
\section{Projects}
    \resumeSubHeadingListStart

    \resumeProjectHeading
      {\textbf{PennOS - UNIX-like Operating System} $|$ \emph{C, POSIX, GDB}}{Mar 2025 -- May 2025}
      \resumeItemListStart
        \resumeItem{Built a user-level operating system in C with a team of 4 (8,000+ lines): process control blocks, PID allocation, parent-child relationships, signal delivery (SIGTERM/SIGSTOP/SIGCONT), and zombie reaping}
        \resumeItem{Implemented a preemptive 3-level priority scheduler with 10ms round-robin quanta and a POSIX-style shell with job control (fg/bg/jobs/kill), I/O redirection, pipelines, and script execution}
      \resumeItemListEnd

    \resumeProjectHeading
      {\textbf{CUDA Path Tracer - GPU PBR Renderer} $|$ \emph{CUDA, C++, GLSL}}{Sep 2025 -- Oct 2025}
      \resumeItemListStart
      \resumeItem{Monte Carlo path tracer with BSDF shading kernels for global illumination, multiple importance sampling, and subsurface scattering; integrated tinyGLTF for glTF 2.0 loading and NVIDIA OptiX for denoising}
      \resumeItem{Profiled with Nsight and optimized via BVH acceleration (3$\times$--160$\times$ framerate on complex scenes), stream compaction (+24--67\%), Russian roulette (+6--24\%), and material sorting. \href{https://github.com/tonytgrt/CUDA-Path-Tracer}{\uline{github.com/tonytgrt/CUDA-Path-Tracer}}}
      \resumeItemListEnd

    \resumeProjectHeading
      {\textbf{MatForge - Vulkan Ray Tracing Material Renderer} $|$ \emph{Vulkan, Slang, C++}}{Nov 2025 -- Dec 2025}
      \resumeItemListStart
      \resumeItem{Implemented four SIGGRAPH papers (2023--2024) in a unified Vulkan ray tracing pipeline, contributing 2,400+ lines of C++ and Slang for sampling and geometry systems}
      \resumeItem{Built a base-3 Sobol' sampler with Owen scrambling across 47 dimensions, achieving 44.7\% MSE reduction vs. PCG sampling at 512 SPP with $<$1\% overhead, plus an automated convergence benchmarking framework}
      \resumeItemListEnd

    \resumeSubHeadingListEnd

\enlargethispage{20pt}

%-----------PROGRAMMING SKILLS-----------
\section{Technical Skills}
 \begin{itemize}[leftmargin=0.15in, label={}]
    \small{\item{
    \textbf{Languages}{: C/C++, Python, Java, Objective-C, Swift, CUDA, C\#, TypeScript, GLSL/MSL/WGSL, Bash} \\
    \textbf{Systems}{: Linux, eBPF, multithreading \& concurrency, POSIX, GDB, CMake, Make, Git, Docker, Xcode} \\
    \textbf{Graphics \& GPU}{: Metal, Vulkan, OpenGL, WebGPU, NVIDIA Nsight, Unreal Engine 5, Unity}
    }}
 \end{itemize}

%-------------------------------------------
\end{document}
